\documentclass[12pt,a4paper]{article}
\usepackage[utf8]{inputenc}
\usepackage[russian]{babel}
\usepackage{geometry}
\geometry{left=2cm,right=2cm,top=2cm,bottom=2cm}

\begin{document}

\section*{Анализ масштабируемости}
В рамках данной лабораторной работы был реализован параллельный алгоритм умножения матрицы на вектор (DGEMV) с применением OpenMP и статического разбиения строк. Анализ полученных результатов подтверждает высокую эффективность распараллеливания при малом количестве вычислительных потоков: до 8 потоков наблюдается практически линейный рост ускорения. Однако при дальнейшем увеличении числа потоков (вплоть до 40) темпы прироста производительности существенно снижаются. При этом показатели ускорения для матриц различной размерности ($M=20\,000$ и $M=40\,000$) демонстрируют схожую динамику.

Отклонение от идеального линейного ускорения на большом числе потоков объясняется тем, что задача DGEMV относится к классу \textit{memory-bound} (ограниченных пропускной способностью памяти). При работе с матрицами большого объёма (до $\sim$12 ГБ) данные перестают помещаться в кэш процессора. В результате потоки начинают конкурировать за общую шину памяти, и итоговая производительность лимитируется скоростью обмена с оперативной памятью, а не вычислительной мощностью ядер.

\end{document}\documentclass[12pt,a4paper]{article}
\usepackage[utf8]{inputenc}
\usepackage[russian]{babel}
\usepackage{geometry}
\geometry{left=2cm,right=2cm,top=2cm,bottom=2cm}

\begin{document}

\section*{Анализ масштабируемости}
В рамках данной лабораторной работы был реализован параллельный алгоритм умножения матрицы на вектор (DGEMV) с применением OpenMP и статического разбиения строк. Анализ полученных результатов подтверждает высокую эффективность распараллеливания при малом количестве вычислительных потоков: до 8 потоков наблюдается практически линейный рост ускорения. Однако при дальнейшем увеличении числа потоков (вплоть до 40) темпы прироста производительности существенно снижаются. При этом показатели ускорения для матриц различной размерности ($M=20\,000$ и $M=40\,000$) демонстрируют схожую динамику.

Отклонение от идеального линейного ускорения на большом числе потоков объясняется тем, что задача DGEMV относится к классу \textit{memory-bound} (ограниченных пропускной способностью памяти). При работе с матрицами большого объёма (до $\sim$12 ГБ) данные перестают помещаться в кэш процессора. В результате потоки начинают конкурировать за общую шину памяти, и итоговая производительность лимитируется скоростью обмена с оперативной памятью, а не вычислительной мощностью ядер.

\end{document}